\begin{helper}[Fixed-base-vertex lifted-size probability tail]
Fix \(n,m\), an edge probability \(p\), a side and a base vertex \(x\), and
let \(0<\delta\le 1/2\).  Assume \(n\le m^2\), \(m>0\), and \(0\le p\le1\).
For samples in which the corresponding red or blue base degree of \(x\) is
within relative error \(\delta\) of \(pm\), the two-bite probability that the
lifted neighborhood of \(x\) is not within relative error \(3\delta\) of
\(pn\) is at most
\[
  \exp\!\left(-(\delta+(1-\delta)\log(1-\delta))(1-\delta)pn\right)
  +
  \exp\!\left(-(((1+\delta)\log(1+\delta)-\delta)(1-\delta)pn)\right).
\]
\end{helper}
\begin{proof}
Condition on the two base graphs.  If the selected base degree of \(x\) is not
within relative error \(\delta\) of \(pm\), the conditioned event is empty.
Otherwise let \(K\subseteq\operatorname{Fin}(m)\times\operatorname{Fin}(m)\)
be the marked set \(N_G(x)\times\operatorname{Fin}(m)\) on the red side, or
\(\operatorname{Fin}(m)\times N_G(x)\) on the blue side.  Its cardinality is
\(|N_G(x)|m\).  The base-degree hypothesis implies
\[
  \operatorname{WithinRelativeError}\!\left(
    \delta,\; pn,\; n|K|/m^2
  \right),
\]
by dividing the degree estimate by \(m\) and multiplying by \(n\).
For a fixed pair of base graphs, the lifted-neighborhood size is exactly the
number of points in the injection image lying in \(K\).  Therefore
\(\noderef{FiberAndDegreeLiftedSizeFixedVertexTail}\) bounds the bad
injections by the displayed Chernoff factor times the total number of
injections.  Multiplying by the uniform injection weight cancels the total
number of injections.  Finally, summing over the two independent base graphs
uses \(\noderef{GnpGraphWeightSumOne}\) twice, while nonnegativity follows
from \(0\le p\le1\).
\end{proof}
